% Part: first-order-logic
% Chapter: syntax-and-semantics
% Section: subformulas

\documentclass[../../../include/open-logic-section]{subfiles}

\begin{document}

\olfileid[fr]{fol}{syn}{sbf}

\olsection{\printtoken{P}{subformula}}

\begin{explain}
Il est souvent utile de parler des !!{formula}s qui « composent »
une !!{formula} donnée. Nous les appelons ses
\emph{!!{subformula}s}. Toute !!{formula} est !!a{subformula}
d'elle-même ; une sous-formule de~$!A$ distincte de~$!A$ est une
\emph{!!{subformula} propre}.
\end{explain}

\begin{defn}[Sous-formules immédiates]
Si $!A$ est !!a{formula}, les \emph{!!{subformula}s immédiates de~$!A$}
sont définies inductivement comme suit :
\begin{enumerate}
\item Une !!{formula} atomique n'a aucune !!{subformula} immédiate.

\tagitem{prvNot}{\indcase{!A}{\lnot !B}{la seule !!{subformula}
    immédiate de $\indfrm$ est~$!B$.}}{}

\item \indcase{!A}{(!B \ast !C)}{les !!{subformula}s immédiates de
  $\indfrm$ sont $!B$ et $!C$ ($\ast$ désigne n'importe quel
  connecteur binaire).}

\tagitem{prvAll}{\indcase{!A}{\lforall[x][!B]}{la seule
    !!{subformula} immédiate de $\indfrm$ est~$!B$.}}{}

\tagitem{prvEx}{\indcase{!A}{\lexists[x][!B]}{la seule
    !!{subformula} immédiate de $\indfrm$ est~$!B$.}}{}
\end{enumerate}
\end{defn}

\begin{defn}[Sous-formules propres]
Si $!A$ est !!a{formula}, les \emph{!!{subformula}s propres de~$!A$}
sont définies récursivement comme suit :
\begin{enumerate}
\item Une !!{formula} atomique n'a aucune !!{subformula} propre.

\tagitem{prvNot}{\indcase{!A}{\lnot !B}{les !!{subformula}s propres
    de $\indfrm$ sont~$!B$ ainsi que toutes les !!{subformula}s
    propres de~$!B$.}}{}

\item \indcase{!A}{(!B \ast !C)}{les !!{subformula}s propres de
  $\indfrm$ sont $!B$ et $!C$, ainsi que toutes les !!{subformula}s
  propres de~$!B$ et de~$!C$.}

\tagitem{prvAll}{\indcase{!A}{\lforall[x][!B]}{les
    !!{subformula}s propres de $\indfrm$ sont~$!B$ ainsi que toutes
    les !!{subformula}s propres de~$!B$.}}{}

\tagitem{prvEx}{\indcase{!A}{\lexists[x][!B]}{les
    !!{subformula}s propres de $\indfrm$ sont~$!B$ ainsi que toutes
    les !!{subformula}s propres de~$!B$.}}{}
\end{enumerate}
\end{defn}

\begin{defn}[Sous-formules]
Les !!{subformula}s de~$!A$ sont $!A$ elle-même et toutes ses
!!{subformula}s propres.
\end{defn}

\begin{explain}
Remarquons la différence subtile entre ces deux définitions. Dans le
premier cas, nous avons directement indiqué les !!{subformula}s
immédiates de~$!A$ pour chacune des formes possibles de~$!A$. Cette
définition explicite par cas reproduit les clauses de la définition
inductive de l'ensemble des !!{formula}s. Dans le second cas, nous
avons aussi suivi ces clauses, mais en ajoutant, pour chaque cas,
les !!{subformula}s propres des !!{formula}s plus petites $!B$ et
$!C$, en plus de $!B$ et~$!C$ elles-mêmes. La définition est donc
\emph{récursive}. En général, une définition d'une fonction sur un
ensemble défini inductivement, ici celui des !!{formula}s, est
récursive si ses clauses font appel à la fonction elle-même. Pour
qu'elle soit bien définie, ces appels ne doivent porter que sur des
arguments qui viennent « avant » celui dont nous définissons la
valeur. Ainsi, pour déterminer les !!{subformula}s propres de
$(!B \ast !C)$, nous utilisons seulement celles des !!{formula}s
antérieures en question.
\end{explain}

\begin{prop}
\ollabel{prop:subfrm-trans}
Si $!B$ est !!a{subformula} de~$!A$ et $!C$ !!a{subformula} de~$!B$,
alors $!C$ est !!a{subformula} de~$!A$. Autrement dit, la relation
« être une sous-formule de » est transitive.
\end{prop}

\begin{prob}
Démontrer la \olref[fol][syn][sbf]{prop:subfrm-trans}.
\end{prob}

\begin{prop}
\ollabel{prop:count-subfrms}
Si $!A$ est une !!{formula} contenant $n$ connecteurs et
quantificateurs, alors $!A$ possède au plus $2n+1$ !!{subformula}s.
\end{prop}

\begin{prob}
Démontrer la \olref[fol][syn][sbf]{prop:count-subfrms}.
\end{prob}

\end{document}
